\section{The crossing-number formula}
\label{sec:formula}

We now combine the preceding results to prove Theorem~\ref{thm:main}.

\subsection{The upper bound}
We recall the classical axis construction; see
Woodall~\cite[p.~658]{Woodall1993}.
Place the vertices $a_i$ on the vertical axis, divided as evenly as
possible between its positive and negative rays. Place the vertices $b_p$
on the horizontal axis in the same way, and join every required pair by
a straight segment. All positions are distinct and nonzero.
An edge interior is contained in one open quadrant.

In a quadrant whose bounding rays contain $u$ first-class and $v$
second-class vertices, each choice of two vertices on each ray forms a
convex quadrilateral. Exactly its diagonals cross, so the contribution is
$\binom u2\binom v2$.
Summing over the quadrants and using \eqref{eq:dr} gives exactly
$d_m d_n$ crossings, provided no three segments concur in their interiors.

The positions can be chosen to ensure this condition. Three segments
concurrent in their interiors must have six distinct endpoints, since
segments with a common endpoint do not cross elsewhere. Their concurrence
is a proper polynomial condition on their axis intercepts: keeping two
segments fixed and varying an intercept of the third avoids their
intersection. There are only finitely many triples, and the product of
their nonzero defining polynomials cannot vanish on the open set of
permitted positions. Thus an admissible drawing achieves the stated count.
In particular,
\begin{equation}
\label{eq:upper}
 \crn(K_{m,n})\le d_m d_n.
\end{equation}
If $m\le2$ or $n\le2$, this already proves equality by nonnegativity.

\subsection{The lower bound for odd \texorpdfstring{$n$}{n}}
Assume $m\ge3$ and $n=2s+1\ge3$, and put $d=s^2=d_n$.
Choose a minimum drawing $D$ normalized as in
Lemma~\ref{lem:normalization}. In particular, no two edges with a common
endpoint cross. Form the blocks $J^{ik}$ of its signed-intersection matrix.
Use the map $\Phi$ from Lemma~\ref{lem:detector}, with any one tuple of
distinct nodes, and set
\[
 R_{ik}=\Phi(J^{ik}),\qquad
 X_i=R_{i1},\qquad Y_k=R_{11}-R_{1k}.
\]
Lemma~\ref{lem:additive} and the annihilation property of $\Phi$ imply
\begin{equation}
\label{eq:rank-family}
 R_{ik}=X_i-Y_k.
\end{equation}
Thus the matrices have the difference form required by
Corollary~\ref{cor:rank-sum}.  We next check its invertibility hypothesis
on diagonal blocks and bound the other block ranks by crossings.

All edges in the star at $a_i$ are oriented outward and do not cross one
another. Consequently $J^{ii}$ is exactly half the sign matrix of the
linear spoke order at $a_i$. Lemma~\ref{lem:detector} gives
\begin{equation}
\label{eq:diagonal-ranks}
 \rank R_{ii}=d\qquad(1\le i\le m).
\end{equation}

For $i\ne k$, let $c_{ik}$ be the number of crossing points between
edges incident with $a_i$ and edges incident with $a_k$.
The entries $J^{ik}(p,p)$ may contain a local term at $b_p$, but $\Phi$
annihilates all diagonal entries: the factor $x_p-x_q$ in
\eqref{eq:detector-definition} makes $B_{pp}=0$.
If $p\ne q$, the two edges $ip,kq$
have no common endpoint, so $J^{ik}(p,q)=I(ip,kq)$.
Each nonzero such entry requires at least one crossing and contributes
to $R_{ik}$ a matrix of rank at most one. Rank subadditivity yields
\begin{equation}
\label{eq:crossing-rank}
 \rank R_{ik}\le c_{ik}\qquad(i\ne k).
\end{equation}
This estimate counts crossing points correctly even when a pair of
edges crosses repeatedly: cancellation in the signed sum can only reduce
the number of nonzero contributions.

Corollary~\ref{cor:rank-sum}, applied to
\eqref{eq:rank-family}--\eqref{eq:diagonal-ranks}, now gives
\[
 2d_m d
 \le \sum_{i\ne k}\rank R_{ik}
 \le \sum_{i\ne k}c_{ik}
 =2c(D).
\]
The last equality holds because each crossing belongs to exactly two
ordered pairs of distinct first-class stars.
Thus $\crn(K_{m,n})\ge d_m d_n$ for odd $n\ge3$.

\subsection{The even case}
We use the vertex-deletion count appearing in
Kleitman~\cite[Section~3, equations~(6)--(7)]{Kleitman1970} and
Woodall~\cite[Theorem~1 and Corollary~1.1]{Woodall1993}.
Let $n>2$ be even, and start again with a normalized minimum drawing $D$.
Delete each of its $n$ second-class vertices in turn. Each remaining
$K_{m,n-1}$ drawing has at least $d_m d_{n-1}$ crossings by the odd case.
Since the two edges at any crossing have distinct second-class endpoints,
that crossing remains in exactly $n-2$ of the deleted-vertex drawings.
Therefore
\[
 (n-2)c(D)\ge n d_m d_{n-1}=(n-2)d_m d_n,
\]
where the last identity follows directly from \eqref{eq:dr} for even $n$.
Division by $n-2$ and the upper bound \eqref{eq:upper} prove equality.
Together with the zero-factor cases, this completes
Theorem~\ref{thm:main}. \qed
